\section{The normalizing algebra and atomless kernels}\label{sec:kernels}

For the rest of the structural argument, $M$ has separable predual,
$A\subseteq M$ is a masa, and $\varphi$ is a faithful $A$-central normal
state.  Lemma~\ref{lem:separable-reduction} permits these assumptions.
We separate the algebra generated by normalizers from its orthogonal
complement with respect to a state-preserving expectation.  The former
has a measured-relation description.  On the latter, we prove that the
positive maps $f\mapsto E_A(y^*fy)$ have atomless kernels.  This
atomlessness will make fine-partition collision terms vanish in the
moment comparison.

\subsection{Normalizers and the measured relation}

A \emph{groupoid normalizer} of $A$ is a partial isometry $v\in M$ such
that $v^*v,vv^*\in A$ and
\[
 vAv^*=Avv^*,\qquad v^*Av=Av^*v.
\]
Let $N$ be the von Neumann algebra they generate.  For every such $v$,
the modular transform $\sigma_t^\varphi(v)$ induces the same
isomorphism between the two corners of $A$, since $\sigma_t^\varphi$
fixes $A$ pointwise.  Consequently
$v^*\sigma_t^\varphi(v)$ is a unitary in $Av^*v$.
These unitaries form a strongly continuous one-parameter group:
the cocycle identity becomes the group identity because the modular
group fixes $Av^*v$.  Thus there is a positive nonsingular operator
$h_v$ affiliated with $Av^*v$ such that
\begin{equation}\label{eq:normalizer-modular}
 \sigma_t^\varphi(v)=v h_v^{it}\qquad(t\in\mathbb R).
\end{equation}
In particular $N$ is invariant under $\sigma^\varphi$.  Takesaki's
Theorem \cite{Tak72} gives a faithful normal $\varphi$-preserving
conditional expectation
\[
 E_N\colon M\longrightarrow N,
 \qquad M^\circ=\ker E_N.
\]
The normalizers
$v1_{[-k,k]}(\log h_v)$ are entire analytic for $\sigma^\varphi$ and
converge strongly-* to $v$.  Indeed Equation~\eqref{eq:normalizer-modular}
extends on this cut to $v h_v^{iz}1_{[-k,k]}(\log h_v)$ for
$z\in\mathbb C$, a norm-entire function.

Identify $A=L^\infty(X,\mu)$ on a standard probability space, with
$\varphi|_A$ given by integration.  A groupoid normalizer induces a
measure-class preserving isomorphism between measurable subsets of
$X$.  Its restriction to the fixed points belongs to $A$.  The remaining
domain can be partitioned into countably many measurable pieces each
disjoint from its image: use a countable family separating distinct
points to obtain a countable cover by such pieces, and then disjointize
the cover.  On one piece let the cut normalizer be $v'$.  Its initial
and final projections are orthogonal, and
\[
 v'+{v'}^*+(1-{v'}^*v'-v'{v'}^*)
\]
is a unitary normalizer.  Recovering $v'$ by a corner cut and summing
strongly shows that full unitary normalizers also generate $N$.
The pair $A\subseteq N$ is therefore Cartan: $A$ is maximal abelian,
regular, and the range of the faithful normal expectation $E_A|_N$.

The Feldman--Moore Representation Theorem
\cite[Theorem~1]{FM77} realizes this pair as the algebra of a countable
nonsingular measured equivalence relation $\mathcal R$ on $(X,\mu)$,
possibly twisted by a scalar $2$-cocycle.  The separably acting
hypothesis holds in the faithful state GNS representation.  We use the
right counting representation
\begin{equation}\label{eq:right-counting-space}
 L^2(\mathcal R,\nu_{\rm r})
 =\int_X^\oplus \ell^2([v]_{\mathcal R})\,d\mu(v),
 \qquad
 \int f\,d\nu_{\rm r}
 =\int_X\sum_{u\sim v}f(u,v)\,d\mu(v).
\end{equation}
Here $A$ acts diagonally by the function of the left site.  With a
normalized cocycle $s$, a kernel $a$ acts in the fiber rooted at $v$ by
the matrix
\[
 (a(u,w)s(u,w,v))_{u,w\in[v]_{\mathcal R}}.
\]
The identity
\[
 s(u,w,v)s(u,v,v')
 =s(u,w,v')s(w,v,v')
\]
shows that rerooting identifies these matrices by diagonal unitary
conjugacy.  Thus their norms and their conjugacy-invariant principal
matrix data do not depend on the chosen root.

Partial isomorphisms with graphs in $\mathcal R$ lift to kernel
groupoid normalizers.  Their finite sums with bounded base
coefficients form a weakly dense *-algebra, denoted $N_{\rm alg}$;
see \cite[Section~2]{FM77}.  The canonical expectation to $A$ takes the
diagonal.  It agrees with $E_A|_N$: on the fixed part of a normalizer
both maps retain its diagonal value, and on pieces with disjoint
initial and final supports both vanish by $A$-bimodularity.  The
identity extends by normality.  Therefore $\varphi|_N$ is diagonal
integration, represented in Equation~\eqref{eq:right-counting-space}
by the unit site vectors $\delta_v$.
All data are understood modulo null sets, with Borel versions on
standard spaces.  Countably many null exceptions may be discarded
along with their relation saturation, which remains null by
nonsingularity.

\subsection{A type I observation}

The kernel proof will produce an intertwiner whose initial projection
is not yet known to belong to $A$.  The following observation is the
step that turns this intertwiner into a groupoid normalizer.  Its state
hypothesis excludes the diffuse masas that can occur in a type I
algebra without a normal conditional expectation.  The importance of
that expectation hypothesis is also reflected in the type I criterion
of \cite[Theorem~4.1]{AS12}; we prove the precise normalizer and
abelian-projection consequences needed here.

\begin{lemma}\label{lem:type-i-normalizers}
Let $D$ be a type I von Neumann algebra with separable predual, and let
$C\subseteq D$ be a masa.  Suppose $D$ admits a faithful $C$-central
normal state.  Then the groupoid normalizers of $C$ generate $D$.
Every abelian projection of $D$ is Murray--von Neumann equivalent in
$D$ to a projection of $C$.
\end{lemma}

\begin{proof}
Use the standard direct integral decomposition of a type I algebra
over its center into full operator algebras; see
\cite[Sections~6.6 and 14.1--14.2]{KR86}.
Integrating ordinary fiber traces against a probability measure in the
center measure class gives a faithful normal semifinite trace $\tau$.
The given state has a positive trace density $h$ of full support, with
$\tau(h)=1$; this is the predual identification in the semifinite case
\cite[Proposition~1.5 and Theorem~2.1]{Haa79}.
Every unitary of $C$ preserves the state and hence its
density.  Thus $h$ commutes with $C$ in the affiliated-operator sense,
and its spectral projections belong to $C$ by maximal abelianness.
The increasing projections
\[
 f_n=1_{[1/n,\infty)}(h),\qquad n\geq1,
\]
exhaust $1$ and satisfy $\tau(f_n)\leq n$.  Their successive
differences form a countable partition of $1$ in $C$ into trace-finite
projections.  Such a projection has finite rank almost everywhere in
the type I fibers.  Split it further by central projections according
to its rank.  Each resulting corner is a measurable field of fixed
finite matrix size, with its compressed $C$ a masa.

We justify explicitly the measurable diagonalization needed in these
finite corners.  Trivialize the finite-rank field and choose countably
many commuting self-adjoint measurable matrix fields generating the
compressed $C$.  In each fiber, the algebras generated by the first
$j$ fields eventually equal the algebra generated by the whole list,
since their dimensions are nondecreasing bounded integers.  The first
index attaining the limiting dimension is measurable: each finite
dimension is determined by ranks of finite lists of polynomial
matrices, and the limiting dimension is their countable supremum.
Split the center according to this first index.  On each piece,
successive eigenspace decompositions measurably diagonalize the finite
commuting list.  Ordered eigenvalues give measurable spectral
projections, and Gram--Schmidt applied to fixed coordinate vectors
gives measurable orthonormal bases of their ranges.  The resulting
rank-one projection fields commute with all generators, not just the
chosen initial list.  As bounded elements of the ambient field
algebra, they therefore belong to the masa.

We have obtained countably many rank-one fields $e_j\in C$, allowing
zero outside their central domains, with $\sum_j e_j=1$.
On $e_jDe_j$ the algebra $Ce_j$ consists precisely of the center
scalars over that domain, since the center of $D$ is contained in
$C$.  Measurable matrix units between $e_j$ and $e_k$ on their common
central domain therefore partially normalize $C$.  These matrix
units and center coefficients generate $D$: compress a bounded field
first by finite sums of the $e_j$, express its matrix entries using
these units, and pass to the strong limit.

Finally, an abelian projection $e\in D$ has rank one almost everywhere
on its central support $z$.  Choose, fiber by fiber on $z$, the first
nonzero $e_j$ from the countable decomposition.  This gives a
projection $e'\in C$ with the same central support and rank.  Measurable
unit vectors spanning the ranges of $e$ and $e'$ give a partial
isometry $w\in D$ with $ww^*=e$ and $w^*w=e'$.  The assertion also
holds for $e=0$.
\end{proof}

\subsection{Atomless kernels off the normalizing algebra}

A positive kernel on $X$ means a measurable family of finite positive
Borel measures $K(v,du)$.  It represents a positive normal map
$K\colon A\to A$ if
\[
 (Kf)(v)=\int_X f(u)\,K(v,du)
 \qquad(f\in L^\infty(X,\mu)),
\]
where for each fixed $\mu$-null Borel set $S$, $K(v,S)=0$ for almost
every $v$.  This condition makes the formula well defined on
equivalence classes; it does not assert that every row measure is
absolutely continuous with respect to $\mu$.

\begin{proposition}\label{prop:atomless-kernels}
In the setting of this section, for every $y\in M^\circ$ the map
\[
 K_y(f)=E_A(y^*fy),\qquad f\in A,
\]
has a positive kernel whose mass is at most $\|y\|^2$ and which is
atomless for almost every output point $v\in X$.
For $k\geq1$ and $y_1,\ldots,y_k\in M^\circ$, the composition
$K_{y_1}\circ\cdots\circ K_{y_k}$ has a positive kernel whose mass is at most
$\prod_{j=1}^k\|y_j\|^2$ and which is atomless almost everywhere.
If $p_{\rm at}\in A$ is the projection of
the atomic part of $A$, then
\[
 p_{\rm at}y=yp_{\rm at}=0\qquad(y\in M^\circ).
\]
\end{proposition}

\begin{proof}
Let $\Omega$ be the standard state vector for $\varphi$.
The commuting left and right actions of $A$ have a joint spectral
measure on $X\times X$ on the vector $y\Omega$.
Since $A\subseteq M_\varphi$, the right action of $f\in A$ on this
vector is $yf\Omega$.  Thus the joint vector measure $\lambda$ satisfies
\begin{equation}\label{eq:joint-kernel-measure}
 \lambda(S\times T)
 =\varphi(y^*1_Sy1_T)
 =\int_T E_A(y^*1_Sy)\,d\mu
\end{equation}
for Borel $S,T\subseteq X$.
Disintegrating relative to the right coordinate gives a kernel for
$K_y$ with mass $E_A(y^*y)(v)\leq\|y\|^2$.
Choose a Borel version and set its rows to zero on a Borel null exceptional
set so that this mass bound holds for every row.
For every fixed $\mu$-null $S$, Equation~\eqref{eq:joint-kernel-measure}
gives $K_y(v,S)=0$ almost everywhere.  This construction uses the joint
spectral measure and does not require product-measure absolute
continuity.

Suppose this kernel is not almost everywhere atomless.  The function
$(v,u)\mapsto K_y(v,\{u\})$ is measurable: express it as the decreasing
limit of the masses of the pieces containing $u$ in an increasing sequence
of finite Borel partitions that separates points. For some integer $k\geq1$, the Borel
set where this function is at least $1/k$ has finite vertical sections
and a projection of positive measure.  The Lusin--Novikov Theorem
\cite[Theorem~18.10]{Kec95} supplies a measurable selection
$s\colon T\to X$ on a positive-measure Borel set $T$ such that
\begin{equation}\label{eq:selected-atom}
 K_y(v,\{s(v)\})\geq1/k\qquad(v\in T).
\end{equation}
For a fixed $\mu$-null Borel set $S$, this implies
\[
 \mu(\{v\in T:s(v)\in S\})
 \leq k\int_T K_y(v,S)\,d\mu(v)=0.
\]
Consequently
\[
 \theta(f)=1_T(f\circ s)\colon A\longrightarrow A1_T
\]
is a well-defined unital *-homomorphism into the corner.  It is normal:
the pushforward under $s$ of any finite measure with an $L^1$ density
on $T$ is absolutely continuous with respect to $\mu$, and its
Radon--Nikodym density gives the preadjoint of $\theta$.

Choose increasing finite Borel partitions $P_j$ of $X$ which generate
$A$ and separate points, and put
\begin{equation}\label{eq:graph-intertwiner-cuts}
 b_j=\sum_{e\in P_j} e y\theta(e).
\end{equation}
The left support projections and the right support projections are
orthogonal families, so $\|b_j\|\leq\|y\|$.
Also $E_N(b_j)=0$ by $N$-bimodularity of $E_N$.
Their state vectors are the decreasing joint projection cuts of
$y\Omega$ to the sets where the left coordinate and $s$ of the right
coordinate lie in the same piece of $P_j$, with the right coordinate
in $T$.  Their limit $\eta$ therefore satisfies
\[
 \|\eta\|^2
 =\int_T K_y(v,\{s(v)\})\,d\mu(v)>0.
\]
A weak-star convergent subnet of the bounded $b_j$ has a limit
$b\in M^\circ$ with $b\Omega=\eta$.  Indeed testing the vectors
against $z\Omega$, $z\in M$, gives the normal functionals
$\varphi(z^*b_j)$, so their vector and weak-star limits agree.
In particular $b\ne0$.  We have $b=b1_T$ and
\begin{equation}\label{eq:graph-intertwiner}
 fb=b\theta(f)\qquad(f\in A).
\end{equation}
For partition step functions this follows from refinement in
Equation~\eqref{eq:graph-intertwiner-cuts}; bounded strong approximation
and normality of $\theta$ extend it to all $f\in A$.

We show that Equation~\eqref{eq:graph-intertwiner} forces $b\in N$.
Write its polar decomposition as $b=u|b|$.
The intertwining relations imply
\[
 |b^*|\in A,\qquad fu=u\theta(f)\quad(f\in A),
 \qquad e:=u^*u\in D:=\theta(A)'\cap1_T M1_T.
\]
Indeed $bb^*$ commutes with $A$, while $b^*b$ commutes with
$\theta(A)$, and polar decomposition preserves the resulting
intertwining relation.  For $d\in eDe$, the element $udu^*$ commutes
with $A$ and belongs to $M$, so lies in $A$.  It follows that $e$ is an
abelian projection in $D$.
Let $z$ be its central support in $D$.
The abelian-projection characterization of type I algebras
\cite[Sections~6.4--6.6]{KR86} shows that $Dz$ is type I: every nonzero central part of
$Dz$ meets $e$ and hence contains a nonzero abelian projection.
Moreover $A1_T$ is a masa in $D$, so its center, and in particular
$z$, belongs to $A1_T$.

The restricted state on $Dz$, after normalization, is faithful,
normal, and central on the masa $Az$.  By
Lemma~\ref{lem:type-i-normalizers}, the algebra $Dz$ is generated by
groupoid normalizers of $Az$.  These also normalize $A$ partially:
their supports lie in $Az$, and every element of $A$ restricts on $z$
to an element of $Az$.  Therefore $Dz\subseteq N$.
The same lemma supplies $w\in Dz$ and $e'\in Az$ with
\[
 ww^*=e,\qquad w^*w=e'.
\]
Set $v=uw$.  Its range projection $uu^*=\supp|b^*|$ and its initial
projection $e'$ belong to $A$.  Since $w$ commutes with $\theta(A)$,
\[
 v^*Av=w^*u^*Auw=\theta(A)e'\subseteq Ae'.
\]
The left side is a masa in $e'Me'$, by conjugating the masa on the
range corner with $v$.  Hence the inclusion is equality, and $v$ is
a groupoid normalizer of $A$.  We conclude
\[
 u=vw^*\in N,\qquad b=|b^*|u\in N,
\]
contrary to $0\ne b\in\ker E_N$.  The kernel is therefore atomless
almost everywhere.

For the composition assertion, let $K_1,K_2$ be two of the kernels.
Choose one fixed $\mu$-null Borel set $S_0$ outside which every row of $K_2$
is atomless.  For almost every $v$, $K_1(v,S_0)=0$.  For each such $v$
the measure
\[
 (K_1\circ K_2)(v,E)=\int_X K_2(u,E)\,K_1(v,du)
\]
vanishes on every singleton $E$; the conclusion holds for all
singletons simultaneously because the integrand row measures are
atomless off $S_0$.  The mass bound and the fixed-null-set property
also pass to this composition.  Iteration proves the assertion for
any finite composition.

Finally the atomic part of the standard probability space can be
represented by a countable set of positive-mass points.  Atomlessness
gives $E_A(y^*p_{\rm at}y)=0$; faithfulness then gives
$p_{\rm at}y=0$.  Applying the same argument to $y^*\in M^\circ$
gives $yp_{\rm at}=0$.
\end{proof}
